\subsection{Instrumento del cuestionario Kano}

El cuestionario se aplica sobre ocho características del centro de salud tipo A. Para cada característica se formulan dos preguntas: una funcional, que indaga la reacción del entrevistado cuando la característica está presente, y una disfuncional, que indaga la reacción cuando está ausente. Las cinco respuestas permitidas son: \emph{me gusta}, \emph{debe ser así}, \emph{me da igual}, \emph{lo tolero} y \emph{no me gusta}.

\begin{table}[ht]
\centering
\footnotesize
\caption{Parejas de preguntas funcionales y disfuncionales del cuestionario Kano.}
\label{tab:kano_preguntas}
\resizebox{\linewidth}{!}{%
\begin{tabular}{@{}p{1.0cm}p{2.8cm}p{4.6cm}p{4.6cm}@{}}
\toprule
ID & Característica & Pregunta funcional & Pregunta disfuncional \\
\midrule
K1 & Identificación única del paciente & ¿Qué sentiría si el sistema asignara un identificador único al paciente? & ¿Qué sentiría si el sistema no asignara un identificador único? \\
K2 & Agenda por especialidad & ¿Qué sentiría si pudiera agendar citas filtrando por especialidad? & ¿Qué sentiría si solo existiera una lista general de citas? \\
K3 & Recordatorio automático & ¿Qué sentiría si el sistema enviara un recordatorio automático de la cita? & ¿Qué sentiría si no existiera ningún recordatorio? \\
K4 & Triaje por nivel de riesgo & ¿Qué sentiría si el sistema clasificara el riesgo del paciente en el triaje? & ¿Qué sentiría si el triaje no clasificara el riesgo? \\
K5 & Historia clínica electrónica & ¿Qué sentiría si la historia clínica estuviera disponible en el sistema? & ¿Qué sentiría si la historia clínica se mantuviera solo en papel? \\
K6 & Disponibilidad de cupos & ¿Qué sentiría si pudiera consultar los cupos libres en tiempo real? & ¿Qué sentiría si los cupos solo se conocieran al llamar al centro? \\
K7 & Registro de signos vitales & ¿Qué sentiría si el sistema permitiera registrar signos vitales durante el triaje? & ¿Qué sentiría si los signos vitales no se registraran en el sistema? \\
K8 & Estadísticas de atención & ¿Qué sentiría si el sistema mostrara estadísticas diarias de atención? & ¿Qué sentiría si no hubiera estadísticas de atención? \\
\bottomrule
\end{tabular}%
}
\end{table}

La tabla de evaluación utilizada para asignar la categoría de cada respuesta se muestra en la Tabla~\ref{tab:kano_matriz}. Las categorías son: M (básica), O (desempeño), A (atractiva), I (indiferente), R (inversa) y Q (cuestionable).

\begin{table}[ht]
\centering
\footnotesize
\caption{Tabla de evaluación del modelo Kano.}
\label{tab:kano_matriz}
\resizebox{\linewidth}{!}{%
\begin{tabular}{@{}p{2.4cm}p{2.0cm}p{2.0cm}p{2.0cm}p{2.0cm}p{2.0cm}@{}}
\toprule
Funcional $\backslash$ Disfuncional & Me gusta & Debe ser & Da igual & Lo tolero & No gusta \\
\midrule
Me gusta   & Q & A & A & A & O \\
Debe ser   & Q & I & I & I & M \\
Da igual   & Q & I & I & I & M \\
Lo tolero  & Q & I & I & I & M \\
No gusta   & R & R & R & R & Q \\
\bottomrule
\end{tabular}%
}
\end{table}

La categoría final de cada característica se obtiene por mayoría simple de las categorías asignadas a las respuestas de los participantes; en caso de empate se prioriza M sobre O, A, I, R y Q.
